\documentclass[10pt,a4paper]{amsart}
\usepackage[margin=19mm]{geometry}
\usepackage{amsmath,amssymb,mathtools}
\usepackage[scr=rsfs,cal=euler]{mathalfa}
\usepackage{xfrac}
\usepackage[unicode,hidelinks,bookmarks=false]{hyperref}
\newcommand{\ie}{i.e.\ }
\newcommand{\cf}{cf.\ }
\newcommand{\resp}{respectively\ }
\newcommand{\Subsection}{\subsection}
\providecommand{\nobreakdash}{\nobreak\hbox{-}}
\setlength{\emergencystretch}{2em}
\multlinegap0pt
\usepackage{amssymb}
\usepackage{algorithm}\usepackage{algpseudocode}
\newtheorem{lemma}{Lemma}
\newtheorem{theorem}{Theorem}
\title{\textbf{Full- and low-rank exponential Euler integrators for the Lindblad equation}}
\author{Hao Chen, Alfio Borz\`i, Denis Jankovi\'c, Jean-Gabriel Hartmann, Paul-Antoine Hervieux}
\date{Source: arXiv:2408.13601v1 (CC BY 4.0)}
\begin{document}
\maketitle

\noindent\emph{Source and licence.} \textbf{Full- and low-rank exponential Euler integrators for the Lindblad equation}. Version arXiv:2408.13601v1 (CC BY 4.0), distributed by arXiv under the Creative Commons Attribution 4.0 International licence, http://creativecommons.org/licenses/by/4.0/, as recorded in the arXiv licence metadata for this exact version and shown on the abstract page https://arxiv.org/abs/2408.13601v1. Full licence notice: \texttt{licenses/arxiv-2408.13601-CC-BY-4.0.txt}.\par\smallskip

\noindent\textit{Editorial scope.} Counted unit: the article's theorem thm4.1 (source0.tex lines 359--364). The article's complete original proof of it is delivered with it (source0.tex lines 365--424, one \texttt{proof} environment of 60 lines). No mathematics is written, repaired or re-derived: the statement and the proof are the article's own lines, in the article's own notation, and no retained line is substituted or re-flowed.  Retained uncounted as the source-local context the proof needs: lines 88--90; lines 135--168; lines 310--315; lines 330--332.  Nothing was omitted from any retained span, so the package carries no omission record.  No retained line contains a citation, so the wrapper carries no bibliography and no bibliography entry.  No retained line contains a figure environment, an image-inclusion command or drawing code, so no image is delivered and the package declares no asset dependency.  Editorial formatting only, in addition: an AMS \texttt{amsart} preamble that restates the article's own declarations and macros used by the retained lines, namely the environments \texttt{lemma}, \texttt{theorem} on separate counters, together with the standard packages those lines need.  The retained environments print in this document as Lemma 1, Theorem 1 in order of appearance, whereas the article prints the same items with the numbering of its own full document.  The complete native source of the article as distributed in the arXiv e-print for this exact version is bundled unchanged as \texttt{sources/164221/source0.tex}.
\begin{equation}\label{equ1.1}
  \dot{\rho}(t)=-i[H,\rho(t)]+\sum_{k=1}^K\gamma_k\left(L_k\rho(t) L_k^{\dag}-\frac{1}{2}\left\{L_k^{\dag}L_k,\rho(t)\right\}\right),
\end{equation}

\section{An exponential Euler integrator}
In this section, we develop an exponential integrator to solve the Lindblad equation \eqref{equ1.1} on the time interval $[0,T]$, for some $T>0$, equipped with the initial condition
$\rho(0)=\rho_0$, where $\rho_0$ is a Hermitian and positive semidefinite matrix with unit trace.
First, let us define
\begin{equation}\label{equ2.1}
  A:=- i \, H-\frac{1}{2}\sum_{k=1}^K\gamma_k \, L_k^{\dag} \, L_k.
\end{equation}
Then, we can rewrite \eqref{equ1.1} as follows:
\begin{equation}\label{equ2.2}
  \dot{\rho}=A \, \rho+\rho \,A^{\dag}+\sum_{k=1}^K\gamma_k \, L_k \, \rho \,L_k^{\dag}.
\end{equation}
Given a positive integer $N$, we divide the time interval by $N$
subintervals of size $\tau=\frac{T}{N}$, with endpoints
$t_n=n\tau$, $n=0,1,\ldots,N$. Then, integrating the equation \eqref{equ2.2} from $t_n$ to $t_{n+1}$ and using the variation-of-constants formula yields
\begin{equation}\label{equ2.3}
  \rho(t_{n+1})=e^{\tau A}\rho(t_n)e^{\tau A^{\dag}}+\sum_{k=1}^K\gamma_k\int_0^{\tau}e^{(\tau-s)A}L_k \, \rho(t_n+s) \, L_k^{\dag}e^{(\tau-s)A^{\dag}}ds.
\end{equation}
Applying this formula again for $\rho(t_n+s)$ within the integral, we obtain
\begin{eqnarray}\label{equ2.4}
  \rho(t_{n+1})=e^{\tau A} \rho(t_n) \, e^{\tau A^{\dag}}+\sum_{k=1}^K\gamma_k
  \int_0^{\tau}e^{(\tau-s)A}L_k \, e^{sA} \, \rho(t_n) \, e^{sA^{\dag}} \, L_k^{\dag} \, e^{(\tau-s)A^{\dag}}ds+R_{n,1},
\end{eqnarray}
where
\begin{equation}\label{equ2.4b}
  R_{n,1}=\sum_{k=1}^K\sum_{j=1}^K\gamma_k\gamma_j\int_0^{\tau}e^{(\tau-s)A}L_k
  \left(\int_0^se^{(s-v)A}L_j \, \rho(t_n+v) \, L_j^{\dag} \, e^{(s-v)A^{\dag}}dv\right)
  L_k^{\dag} \, e^{(\tau-s)A^{\dag}}ds.
\end{equation}
By truncating the term $R_{n,1}$ and approximating the matrix exponentials by $e^{(\tau-s)A}\approx I$ and $e^{(\tau-s)A^{\dag}}\approx I$ in \eqref{equ2.4}, we obtain our full-rank exponential Euler (FREE) integrator, for the approximation $\{\rho_n\}_{n=0}^N$ to $\{\rho(t_n)\}_{n=0}^N$, as follows:
\begin{equation}\label{equ2.5}
  \rho_{n+1}=e^{\tau A} \rho_n \, e^{\tau A^{\dag}}+\sum_{k=1}^K\gamma_k
  \int_0^{\tau}L_k \, e^{sA} \, \rho_n \, e^{sA^{\dag}} L_k^{\dag}\, ds ,
\end{equation}
for $0\leq n\leq N-1$.

\begin{lemma}\label{lem4.1}
   For any Hermitian matrix $\varrho\in \mathbb{C}^{m\times m}$, it holds that
  \begin{equation*}
    \left\|e^{tA} \, \varrho \, e^{tA^{\dag}}\right\|_1\leq \left\|\varrho\right\|_1,~~~\, t\geq0.
  \end{equation*}
\end{lemma}

\begin{lemma}\label{lem4.3}
  For any Hermitian matrix $\varrho\in \mathbb{C}^{m\times m}$ and any $t\geq0$, it holds that
\end{lemma}

\begin{theorem}\label{thm4.1}
  Let $\rho(t)$ be the solution of the Lindblad equation \eqref{equ1.1} with
  $\rho(0)=\rho_0$, and $\{\rho_n\}_{n=0}^N$ be the corresponding numerical solution generated by the full-rank exponential Euler scheme \eqref{equ2.5}. Then, it holds that
  \[\|\rho(t_n)-\rho_n\|_1\leq c_1t_n\tau,~~~0\leq n\leq N,\]
  where the constant $c_1>0$ depends on $C_1,~C_2,~T$, but is independent of $\tau$ and $n$.
\end{theorem}

\begin{proof}
Expanding $e^{(\tau-s)A}$ in a Taylor series with remainder in integral form yields
\begin{equation*}
  e^{(\tau-s)A}=I+\int_{s}^{\tau}A \, e^{(\tau-v)A}dv.
\end{equation*}
Inserting the above equation into \eqref{equ2.4}, we have
\begin{eqnarray}\label{equ4.2}
  \rho(t_{n+1})=e^{\tau A}\rho(t_n) \, e^{\tau A^{\dag}}+\sum_{k=1}^K\gamma_k
  \int_0^{\tau}L_k \, e^{sA} \, \rho(t_n) \, e^{sA^{\dag}} \, L_k^{\dag} \, ds+\sum_{j=1}^4R_{n,j},
\end{eqnarray}
where $R_{n,1}$ is defined in \eqref{equ2.4b}, and
  \begin{align*}
  &R_{n,2}=\sum_{k=1}^K\gamma_k
  \int_0^{\tau}\left(\int_s^{\tau}Ae^{(\tau-v)A}dv\right)L_ke^{sA}\rho(t_n)e^{sA^{\dag}} L_k^{\dag} \, ds,\\
  &R_{n,3}=\sum_{k=1}^K\gamma_k
  \int_0^{\tau}L_ke^{sA}\rho(t_n)e^{sA^{\dag}} L_k^{\dag}\left(\int_s^{\tau}A^{\dag}e^{(\tau-v)A^{\dag}}dv\right)ds,\\
  &R_{n,4}=\sum_{k=1}^K\gamma_k
  \int_0^{\tau}\left(\int_s^{\tau}Ae^{(\tau-v)A}dv\right)L_ke^{sA}\rho(t_n)e^{sA^{\dag}} L_k^{\dag}\left(\int_s^{\tau}A^{\dag}e^{(\tau-v)A^{\dag}}dv\right)ds.
\end{align*}
Denote the global error $E_n:=\rho(t_n)-\rho_n$. Subtracting \eqref{equ2.5} from \eqref{equ4.2}, we obtain the error recursion
\begin{equation}\label{equ4.3}
  E_{n+1}=e^{\tau A}E_ne^{\tau A^{\dag}}+\sum_{k=1}^K\gamma_k
  \int_0^{\tau}L_ke^{sA}  \, E_n  \, e^{sA^{\dag}} L_k^{\dag} \, ds+R_n,
\end{equation}
where $R_n=\sum_{j=1}^4R_{n,j}$.

Using Lemma \ref{lem4.1} and noting that $\|\rho(t)\|_1=1$ for any $t\geq0$, we have
\begin{eqnarray}\label{equ4.4}
\|R_{n,1}\|_1&\leq&\sum_{k,j=1}^K\gamma_k\gamma_j\int_0^{\tau}\left\|e^{(\tau-s)A}L_k
  \left(\int_0^se^{(s-v)A}L_j\rho(t_n+v)L_j^{\dag}e^{(s-v)A^{\dag}}dv\right)
  L_k^{\dag}e^{(\tau-s)A^{\dag}}\right\|_1ds\nonumber\\
  &\leq& \sum_{k,j=1}^K\gamma_k\gamma_j\int_0^{\tau}\|L_k\|_1^2
  \int_0^s\left\|e^{(s-v)A}L_j\rho(t_n+v)L_j^{\dag}e^{(s-v)A^{\dag}}\right\|_1dv\nonumber\\
  &\leq& \sum_{k,j=1}^K\gamma_k\gamma_j\|L_k\|_1^2\|L_j\|_1^2
  \int_0^{\tau}\int_0^s\|\rho(t_n+v)\|_1dvds\nonumber\\
  &=&\frac{1}{2}\tau^2\sum_{k,j=1}^K\gamma_k\gamma_j\|L_k\|_1^2\|L_j\|_1^2=\frac{1}{2}C_1^2\tau^2.
\end{eqnarray}
Similarly, we can obtain
\begin{equation}\label{equ4.5}
  \|R_{n,2}\|_1\leq\frac{1}{2}C_1 \, C_2 \, \tau^2,~~~\|R_{n,3}\|_1\leq\frac{1}{2}C_1\, C_2 \, \tau^2,
\end{equation}
and
\begin{equation}\label{equ4.6}
  \|R_{n,4}\|_1\leq\frac{1}{3}C_1 \, C_2^2 \, \tau^3.
\end{equation}
Combining \eqref{equ4.4}-\eqref{equ4.6}, we have
\begin{equation}\label{equ4.7}
  \|R_n\|_1\leq c_1 \, \tau^2,
\end{equation}
where $c_1=\frac{1}{2}C_1^2+C_1C_2+\frac{1}{3}C_1C_2^2 \, T$.

It follows from $\eqref{equ4.3}$, $\eqref{equ4.7}$, $E_0=0$ and Lemma \ref{lem4.3} that
\begin{eqnarray*}
  \|E_{n+1}\|_1&\leq&\left\|e^{\tau A}E_ne^{\tau A^{\dag}}+\sum_{k=1}^K\gamma_k
  \int_0^{\tau}L_ke^{sA} \, E_n  \, e^{sA^{\dag}} L_k^{\dag} \, ds\right\|_1+\|R_n\|_1\\
  &\leq& \|E_n\|_1+c_1\tau^2\leq \|E_{n-1}\|_1+2c_1\tau^2\\
  &\leq& \|E_0\|_1+c_1(n+1)\tau^2=c_1 \, t_{n+1} \, \tau,
\end{eqnarray*}
which completes the proof.
\end{proof}

\end{document}
